\setlength{\absparsep}{18pt} % ajusta o espaçamento dos parágrafos do resumo
\begin{resumo}

	Este trabalho apresenta o desenvolvimento de uma plataforma \textit{web} para a criação e o compartilhamento de coleções digitais de qualquer natureza, motivada pela rigidez das ferramentas usuais de registro, que tratam cada registro como um conjunto de campos fixos de aparência padronizada. Na plataforma, o acervo organiza-se em categorias, coleções e itens, e a ficha de cada item é definida pelo próprio usuário em um editor visual de \textit{layouts}, com elementos de seis tipos. Completam a proposta a personalização da interface por paletas de cores e tipografia e as funcionalidades sociais de amizade, acompanhamento de coleções, \textit{feed} e notificações, sob a privacidade definida em cada coleção. O desenvolvimento partiu da especificação dos requisitos e do modelo de domínio, seguiu uma arquitetura em camadas, com interface em React, API em NestJS e banco de dados PostgreSQL, e foi verificado por testes automatizados, verificação manual e medições de desempenho. A plataforma foi publicada e está em operação. Os 546 testes automatizados, de unidade e de ponta a ponta, foram aprovados, os vinte e quatro requisitos funcionais foram atendidos e as métricas de desempenho foram cumpridas com folga: com a rede limitada a 10 Mbps, as telas medidas exibiram seu conteúdo principal em menos de meio segundo no percentil 75, e o editor respondeu às interações em até 48 milissegundos. Conclui-se que a plataforma reúne, em um só lugar, a organização de coleções, a liberdade de definir a estrutura e a aparência de cada registro e o compartilhamento com outras pessoas.

	\vspace{\onelineskip}

	\textbf{Palavras-chave}: Coleções digitais. Personalização de interfaces. Editor de \textit{layouts}. Aplicação \textit{web}. Compartilhamento.

\end{resumo}
